\documentclass[11pt]{article}

\usepackage[utf8]{inputenc}
\usepackage[T1]{fontenc}
\usepackage{lmodern}
\usepackage{geometry}
\usepackage{amsmath,amssymb,amsfonts,amsthm,mathtools}
\usepackage{bm}
\usepackage{enumitem}
\usepackage{microtype}
\usepackage{hyperref}
\usepackage{titlesec}
\usepackage{booktabs}
\usepackage{array}
\usepackage{setspace}
\usepackage{mathrsfs}
\usepackage{physics}

\geometry{margin=1in}
\hypersetup{colorlinks=true, linkcolor=black, urlcolor=blue, citecolor=black}
\setlist[enumerate]{topsep=0.4em,itemsep=0.35em}
\setlist[itemize]{topsep=0.3em,itemsep=0.25em}
\titleformat{\section}{\large\bfseries}{\thesection.}{0.5em}{}
\titleformat{\subsection}{\normalsize\bfseries}{\thesubsection.}{0.5em}{}
\setstretch{1.06}

\newcommand{\Aether}{\AE{}ther}
\newcommand{\Aflow}{\AE{}ther-flow}
\newcommand{\TheoryName}{\Aflow{} Interpretation of Relativity}
\newcommand{\TheoryProgram}{\Aether{} / \Aflow{} framework}
\newcommand{\Sub}{\mathcal{A}}
\newcommand{\BridgeMetric}{\mathfrak{g}}

\newtheorem{definition}{Definition}
\newtheorem{lemma}{Lemma}
\newtheorem{proposition}{Proposition}
\newtheorem{theorem}{Theorem}
\newtheorem{remark}{Remark}
\newtheorem{corollary}{Corollary}

\title{\textbf{The \TheoryName{}}\\[0.4em]
\large Derivative-Mixing Branch Quadratic Consistency and Flat-Background Exact-Closure Test}
\author{}
\date{}

\begin{document}

\maketitle

\begin{abstract}
This manuscript carries the derivative-mixing \(S/Q\) branch through the next required test stage: full quadratic auxiliary coefficients, flat-background exact-closure conditions, and the decision of whether a nondegenerate observer-accessible branch actually exists. The scope remains disciplined. The note does not claim a completed substrate derivation of Einsteinian gravity. It analyzes only the first nonminimal branch already selected in the active research plan.

The result is sharper than the branch-selection note alone. The derivative-mixing operator changes the scalar coefficient obstruction, but under the branch's own gradient-health condition the new low-momentum coefficient still admits a positive decomposition. Writing
\[
\beta_I\equiv \widehat M^{-1\,IJ}\mu_{SJ},
\]
the generalized quadratic exact-closure coefficient becomes
\[
\mathcal B_{\rho,2}
=
\bigl(\kappa_S-\rho_I\rho_I\bigr)
+
\bigl(\rho_I-\beta_I\bigr)\bigl(\rho_I-\beta_I\bigr).
\]
Hence a nondegenerate flat-background exact-closure branch is impossible once the scalar-gradient block is required to be positive definite. Exact closure survives only on the tuned semidefinite branch
\[
\rho_I=\beta_I,
\qquad
\kappa_S=\rho_I\rho_I,
\qquad
\mu_S^2=\mu_{SI}\beta_I,
\]
for which the entire scalar correction vanishes identically and the \((\sigma,q^I)\) sector factorizes into an exact square. The derivative-mixing branch therefore evades the original minimal no-go theorem only by collapsing onto a new degenerate branch. It does not supply the sought nondegenerate flat-background exact-closure solution.
\end{abstract}

\section{Introduction}

The active exact-closure line of \TheoryName{} is already fixed by the foundations, dynamics, consistency, relativistic-recovery, and flow-geometry manuscripts. The derivational side of the project then advanced through three further stages: an explicit substrate-kinematics bridge, a coefficient-level linearized consistency pass for the minimal candidate action, and a no-go result showing that the minimal stable-sign branch has no nondegenerate flat-background exact-closure solution.

After that no-go result, the first nonminimal branch selected for testing was the derivative-mixing completion of the \(S/Q\) sector. The branch-selection note established why that branch was the natural next target and showed that the old minimal positive-sum obstruction no longer applied in its original form. What remained open was the actual Phase 3 calculation: derive the full quadratic auxiliary coefficients of that branch and determine whether the new cancellation conditions can be satisfied on a nondegenerate flat-background exact-closure branch.

That is the task of the present manuscript. The answer is now definite. The derivative-mixing branch changes the obstruction, but once its own health condition is imposed, it still does not produce a genuinely nondegenerate flat-background exact-closure branch.

\paragraph{Framework and claim boundary.}
\TheoryProgram{} denotes the broader \Aether{} / \Aflow{} framework built on the claims that the \Aether{} is the underlying four-dimensional substrate of reality, the \Aflow{} is its intrinsic ordered motion, observed three-dimensional space is the local experiential slice of that deeper substrate, S-time is the experienced order of change arising from matter, light, and the \Aflow{}, and observed expansion is the three-dimensional appearance of deeper four-dimensional ordered motion. Within that framework, \TheoryName{} denotes the exact-closure theory in which the effective gravitational dynamics are adopted to be exactly Einsteinian with universal matter coupling. In the active sequence, \emph{adoption} means use of that established relativistic dynamics without claiming substrate derivation, while \emph{derivation} is reserved for a first-principles recovery from explicit substrate variables.

The present note stays inside that boundary. It does not promote the derivative-mixing branch to established theory. It determines its actual quadratic status.

\section{Derivative-Mixing Branch and Background Variables}

The derivative-mixing branch modifies the minimal candidate action only through the auxiliary sector,
\begin{equation}
\mathcal L_{\mathrm{aux}}^{(\rho)}
=
\mathcal L_{\mathrm{aux}}
-\rho_I P^{AB}\partial_A S\,D_BQ^I,
\label{eq:Lauxrho}
\end{equation}
while the bridge metric \(\BridgeMetric_{AB}=G_{AB}-2U_AU_B\), the Einsteinian benchmark, and the universal matter-coupling route through \(S_{\mathrm{matter}}[\Xi^*\BridgeMetric,\psi]\) are all left unchanged.

As in the previous flat-background analysis, use the admissible homogeneous bridge background
\begin{equation}
\bar G_{AB}=\delta_{AB},
\qquad
\bar U^A=\delta^A{}_0,
\qquad
\bar S=\sigma_0X^0,
\qquad
\bar Q^I=0,
\label{eq:background}
\end{equation}
for which
\begin{equation}
\bar\BridgeMetric_{AB}=\eta_{AB}=\mathrm{diag}(-1,1,1,1).
\label{eq:etabridge}
\end{equation}
Write perturbations as
\begin{equation}
G_{AB}=\delta_{AB}+H_{AB},
\qquad
U^A=\bar U^A+V^A,
\qquad
S=\sigma_0X^0+\sigma,
\qquad
Q^I=q^I.
\label{eq:perturbations}
\end{equation}
The unit constraint gives
\begin{equation}
V^0=-\frac{1}{2}H_{00},
\label{eq:V0}
\end{equation}
and the bridge perturbation is
\begin{equation}
h_{00}=-H_{00},
\qquad
h_{0i}=-H_{0i}-2V_i,
\qquad
h_{ij}=H_{ij}.
\label{eq:bridgepert}
\end{equation}
As before, decompose
\begin{equation}
V_i=V_i^T+\partial_i\chi,
\qquad
\partial_iV_T^i=0.
\label{eq:Vdecomp}
\end{equation}

\section{Full Quadratic Auxiliary Lagrangian}

\begin{proposition}[Quadratic auxiliary coefficients of the derivative-mixing branch]
On the flat background \eqref{eq:background}--\eqref{eq:etabridge}, the full quadratic auxiliary Lagrangian of the derivative-mixing branch is
\begin{align}
\mathcal L_{\mathrm{aux},\rho}^{(2)}
=\;&
-\frac{m_S^2}{2}\bigl(\partial_0\sigma-\sigma_0\phi\bigr)^2
-\frac{\kappa_S}{2}\,\partial_i\sigma\,\partial_i\sigma
-\frac{1}{2}\mu_S^2\sigma^2
-\mu_{SI}\sigma q^I
\nonumber\\
&-\rho_I\,\partial_i\sigma\,\partial_i q^I
-\frac{1}{2}q^I\!\left(-\nabla^2\delta_{IJ}+\widehat M_{IJ}\right)\!q^J
\nonumber\\
&-\frac{\kappa_\theta}{2}\left(\nabla^2\chi+\frac{1}{2}\partial_0H\right)^2
-\frac{\kappa_\Omega}{4}\,
\partial_{[i}\!\left(S_{j]}+V^T_{j]}\right)
\partial^{[i}\!\left(S^{j]}+V_T^{j]}\right),
\label{eq:Laux2rho}
\end{align}
with \(h_{00}=-2\phi\).
\end{proposition}

\begin{proof}
The minimal quadratic coefficients were already derived in the coefficient-level linearized-consistency manuscript. The only new term is the quadratic expansion of the derivative-mixing operator,
\[
-\rho_I P^{AB}\partial_A S\,D_BQ^I
=
-\rho_I\,\delta^{ij}\partial_i\sigma\,\partial_j q^I
+\mathcal O(3),
\]
because \(P^{AB}\) reduces to the spatial projector on the chosen branch and \(D_BQ^I\) reduces to \(\partial_Bq^I\) at this order. Adding that term to the previously derived minimal quadratic auxiliary sector yields \eqref{eq:Laux2rho}.
\end{proof}

\begin{remark}
The derivative-mixing completion changes only the scalar \((\sigma,q^I)\) sector. The longitudinal and transverse \(U\)-sector pieces still enter exactly as in the minimal branch.
\end{remark}

\begin{proposition}[Unchanged \(U\)-sector elimination]
On the decaying flat branch with no harmonic zero mode, the \(\chi\) and \(V_T^i\) Euler--Lagrange equations remain
\begin{equation}
\nabla^2\chi=-\frac{1}{2}\partial_0H,
\qquad
V_T^i=-S^i.
\label{eq:Uelimination}
\end{equation}
Consequently the \(\kappa_\theta\) and \(\kappa_\Omega\) sectors contribute no residual flat-background bridge self-energy at quadratic order.
\end{proposition}

\begin{proof}
The derivative-mixing term in \eqref{eq:Laux2rho} depends only on \(\sigma\) and \(q^I\). Therefore the variations with respect to \(\chi\) and \(V_T^i\) are unchanged from the minimal branch, and so are their solutions and their vanishing on-shell quadratic contribution.
\end{proof}

\section{Scalar Reduction and Generalized Exact-Closure Operator}

Define the \(Q\)-sector spatial operator
\begin{equation}
\mathcal A_{IJ}(-\nabla^2)
\equiv
(-\nabla^2)\delta_{IJ}+\widehat M_{IJ}.
\label{eq:Aoperator}
\end{equation}
Then the \(q^I\) Euler--Lagrange equation derived from \eqref{eq:Laux2rho} is
\begin{equation}
\mathcal A_{IJ}(-\nabla^2)\,q^J
=
-\bigl(\mu_{SI}+\rho_I(-\nabla^2)\bigr)\sigma.
\label{eq:qeqrho}
\end{equation}
Hence
\begin{equation}
q^I
=
-\mathcal A^{-1\,IJ}(-\nabla^2)\,
\bigl(\mu_{SJ}+\rho_J(-\nabla^2)\bigr)\sigma.
\label{eq:qsolverho}
\end{equation}

Substituting \eqref{eq:qsolverho} into the scalar sector yields
\begin{equation}
\mathcal L_{\sigma,\mathrm{eff}}^{(2)}
=
-\frac{m_S^2}{2}\bigl(\partial_0\sigma-\sigma_0\phi\bigr)^2
-\frac{1}{2}\sigma\,\mathcal B_{\rho}(-\nabla^2)\,\sigma,
\label{eq:Lsigmaeffrho}
\end{equation}
where
\begin{equation}
\mathcal B_{\rho}(-\nabla^2)
=
\kappa_S(-\nabla^2)
+\mu_S^2
-\bigl(\mu_{SI}+\rho_I(-\nabla^2)\bigr)
\mathcal A^{-1\,IJ}(-\nabla^2)
\bigl(\mu_{SJ}+\rho_J(-\nabla^2)\bigr).
\label{eq:Brhooperator}
\end{equation}

In Fourier space,
\begin{equation}
\mathcal B_{\rho}(k^2)
=
\kappa_Sk^2+\mu_S^2
-\bigl(\mu_{SI}+\rho_Ik^2\bigr)
\bigl(k^2\delta+\widehat M\bigr)^{-1\,IJ}
\bigl(\mu_{SJ}+\rho_Jk^2\bigr).
\label{eq:Brhok}
\end{equation}
Its low-momentum expansion is
\begin{equation}
\mathcal B_{\rho}(k^2)
=
\mathcal B_{\rho,0}
+\mathcal B_{\rho,2}k^2
+\mathcal O(k^4),
\label{eq:Brhoexpand}
\end{equation}
with
\begin{align}
\mathcal B_{\rho,0}
&=
\mu_S^2-\mu_{SI}\widehat M^{-1\,IJ}\mu_{SJ},
\label{eq:Brho0}
\\
\mathcal B_{\rho,2}
&=
\kappa_S
+\mu_{SI}\widehat M^{-2\,IJ}\mu_{SJ}
-2\mu_{SI}\widehat M^{-1\,IJ}\rho_J.
\label{eq:Brho2}
\end{align}

As in the minimal branch, flat-background quadratic exact closure requires
\begin{equation}
\mathcal B_{\rho,0}=0,
\qquad
\mathcal B_{\rho,2}=0.
\label{eq:Brhoconditions}
\end{equation}

\section{Gradient Health Condition and Nondegenerate No-Go Result}

The branch-selection note already identified the spatial-gradient block of the \((\sigma,q^I)\) sector as
\begin{equation}
\mathbb K_{\mathrm{grad}}
=
\begin{pmatrix}
\kappa_S & \rho_J \\
\rho_I & \delta_{IJ}
\end{pmatrix}.
\label{eq:Kgrad}
\end{equation}
The branch's basic health requirement is
\begin{equation}
\mathbb K_{\mathrm{grad}}\succeq 0,
\label{eq:Kpositive}
\end{equation}
which is equivalent to
\begin{equation}
\kappa_S-\rho_I\rho_I \ge 0.
\label{eq:Schur}
\end{equation}

\begin{definition}[Nondegenerate gradient-health branch]
For the present flat-background analysis, a derivative-mixing branch is called nondegenerate if
\begin{equation}
\mathbb K_{\mathrm{grad}}\succ 0,
\label{eq:Kstrict}
\end{equation}
equivalently
\begin{equation}
\kappa_S-\rho_I\rho_I>0.
\label{eq:Schurstrict}
\end{equation}
\end{definition}

Define
\begin{equation}
\beta_I \equiv \widehat M^{-1\,IJ}\mu_{SJ}.
\label{eq:beta}
\end{equation}
Then \eqref{eq:Brho2} may be rewritten as
\begin{equation}
\mathcal B_{\rho,2}
=
\bigl(\kappa_S-\rho_I\rho_I\bigr)
+
\bigl(\rho_I-\beta_I\bigr)\bigl(\rho_I-\beta_I\bigr).
\label{eq:Brho2decomp}
\end{equation}

\begin{lemma}[Positive decomposition of the derivative-mixing coefficient]
If \(\widehat M_{IJ}\) is positive definite and \(\mathbb K_{\mathrm{grad}}\succeq 0\), then
\begin{equation}
\mathcal B_{\rho,2}\ge 0.
\label{eq:Brho2positive}
\end{equation}
Moreover,
\begin{equation}
\mathcal B_{\rho,2}=0
\quad\Longleftrightarrow\quad
\rho_I=\beta_I
\ \text{and}\
\kappa_S=\rho_I\rho_I.
\label{eq:Brho2iff}
\end{equation}
\end{lemma}

\begin{proof}
Equation \eqref{eq:Brho2decomp} is an exact identity. Under \eqref{eq:Schur}, both terms on the right-hand side are nonnegative. Therefore \(\mathcal B_{\rho,2}\ge 0\). The right-hand side vanishes only if both terms vanish separately, which is exactly \eqref{eq:Brho2iff}.
\end{proof}

\begin{theorem}[No nondegenerate flat-background exact-closure branch]
Assume \(\widehat M_{IJ}\) is positive definite and impose the derivative-mixing branch health condition \(\mathbb K_{\mathrm{grad}}\succeq 0\). Then:
\begin{enumerate}
    \item the derivative-mixing branch admits no nondegenerate flat-background exact-closure solution satisfying \(\mathbb K_{\mathrm{grad}}\succ 0\);
    \item exact closure can occur only on the tuned semidefinite branch
    \begin{equation}
    \rho_I=\beta_I=\widehat M^{-1\,IJ}\mu_{SJ},
    \qquad
    \kappa_S=\rho_I\rho_I,
    \qquad
    \mu_S^2=\mu_{SI}\beta_I.
    \label{eq:tunedbranch}
    \end{equation}
\end{enumerate}
\end{theorem}

\begin{proof}
If \(\mathbb K_{\mathrm{grad}}\succ 0\), then \(\kappa_S-\rho_I\rho_I>0\). By \eqref{eq:Brho2decomp}, this makes \(\mathcal B_{\rho,2}>0\) regardless of \(\beta_I\). Therefore the exact-closure condition \(\mathcal B_{\rho,2}=0\) cannot be met on any nondegenerate branch.

If instead one allows \(\mathbb K_{\mathrm{grad}}\succeq 0\), then Lemma 1 shows that \(\mathcal B_{\rho,2}=0\) is equivalent to \(\rho_I=\beta_I\) and \(\kappa_S=\rho_I\rho_I\). Substituting \(\rho_I=\beta_I\) into \eqref{eq:Brho0} gives
\[
\mathcal B_{\rho,0}
=
\mu_S^2-\mu_{SI}\beta_I.
\]
Therefore \(\mathcal B_{\rho,0}=0\) requires the third relation in \eqref{eq:tunedbranch}. This proves both claims.
\end{proof}

\begin{corollary}
The derivative-mixing branch improves the minimal no-go result only in a limited sense: it replaces the old degenerate branch by a new tuned degenerate branch, but it still does not provide a genuinely nondegenerate flat-background exact-closure solution.
\end{corollary}

\section{Exact Degenerate Branch and Perfect-Square Factorization}

The tuned branch \eqref{eq:tunedbranch} has an instructive exact form.

\begin{proposition}[Exact square on the tuned branch]
When \eqref{eq:tunedbranch} holds, the scalar \(Q/S\) sector factorizes as
\begin{equation}
\mathcal L_{\sigma q,\rho}^{(2)}
=
-\frac{m_S^2}{2}\bigl(\partial_0\sigma-\sigma_0\phi\bigr)^2
-\frac{1}{2}
\bigl(q^I+\rho_I\sigma\bigr)
\left(-\nabla^2\delta_{IJ}+\widehat M_{IJ}\right)
\bigl(q^J+\rho_J\sigma\bigr).
\label{eq:perfectsquare}
\end{equation}
Equivalently,
\begin{equation}
\mathcal B_{\rho}(k^2)=0
\label{eq:Brhozero}
\end{equation}
for all \(k^2\) at quadratic order.
\end{proposition}

\begin{proof}
Using \(\rho_I=\beta_I\) implies
\begin{equation}
\mu_{SI}=\widehat M_{IJ}\rho_J.
\label{eq:muMrelation}
\end{equation}
The conditions \(\kappa_S=\rho_I\rho_I\) and \(\mu_S^2=\mu_{SI}\rho_I\) then show that every \(\sigma\)-quadratic coefficient in \eqref{eq:Laux2rho} is exactly the coefficient obtained by expanding \eqref{eq:perfectsquare}. Expanding the square gives the \(q\)-sector quadratic term, the derivative-mixing term, the \(\sigma q\) potential mixing term, the \(\sigma\)-gradient term, and the \(\sigma\)-mass term with the required coefficients. Therefore \eqref{eq:perfectsquare} holds.

In Fourier space, \eqref{eq:muMrelation} implies
\[
\mu_{SI}+\rho_Ik^2
=
\bigl(\widehat M_{IJ}+k^2\delta_{IJ}\bigr)\rho_J.
\]
Substituting this into \eqref{eq:Brhok} yields
\[
\bigl(\mu_{SI}+\rho_Ik^2\bigr)
\bigl(k^2\delta+\widehat M\bigr)^{-1\,IJ}
\bigl(\mu_{SJ}+\rho_Jk^2\bigr)
=
\rho_I\widehat M_{IJ}\rho_J+\rho_I\rho_I\,k^2.
\]
Using the tuned-branch conditions then gives \(\mathcal B_{\rho}(k^2)=0\) exactly.
\end{proof}

\begin{remark}
The tuned branch has exact quadratic closure in the observer-accessible bridge sector, but it is still degenerate. The scalar-gradient block saturates the semidefinite bound, and the independent scalar combination left outside the exact square carries only the Stueckelberg-like time-derivative structure \((\partial_0\sigma-\sigma_0\phi)^2\).
\end{remark}

\begin{corollary}[Observer-accessible pole statement]
On the tuned branch \eqref{eq:tunedbranch}, no unsuppressed observer-accessible scalar, vector, or tensor correction survives in the flat-background bridge self-energy at quadratic order. The tensor sector remains Fierz--Pauli, the \(U\)-sector remains auxiliary, and the scalar bridge self-energy vanishes identically.
\end{corollary}

\section{What This Phase 3 Test Establishes}

The derivative-mixing branch now has a definite quadratic status.
\begin{enumerate}
    \item The full quadratic auxiliary coefficients are known explicitly.
    \item The \(\chi\) and \(V_T^i\) sectors remain auxiliary exactly as in the minimal branch.
    \item The scalar exact-closure burden is compressed to the generalized operator \(\mathcal B_{\rho}(k^2)\).
    \item Under the branch's own gradient-health condition, the low-momentum exact-closure coefficient still admits a positive decomposition.
    \item Therefore no nondegenerate flat-background exact-closure branch exists.
    \item Exact closure survives only on the tuned semidefinite branch \eqref{eq:tunedbranch}, where the scalar sector factorizes into an exact square and the bridge self-energy vanishes identically.
\end{enumerate}

This outcome is scientifically useful even though it is not the hoped-for positive derivational completion. It means the derivative-mixing branch has now been tested hard enough to know its actual role. It is not the nondegenerate exact-closure branch sought by the post-no-go program.

\section{Conclusion}

The Phase 3 task for the derivative-mixing branch was to derive its quadratic auxiliary coefficients and determine whether it actually possesses a nondegenerate flat-background exact-closure branch. That question now has a clear answer.

The derivative-mixing completion does alter the minimal obstruction, but it does not remove the deeper pattern. Once the branch's own gradient-health condition is enforced, the new low-momentum coefficient \(\mathcal B_{\rho,2}\) again becomes a sum of nonnegative pieces. As a result, nondegenerate exact closure is impossible. The only exact-closure solution is the tuned semidefinite branch \eqref{eq:tunedbranch}, where the scalar/Q sector collapses into an exact square and the bridge self-energy vanishes identically.

The derivative-mixing branch therefore earns a narrower scientific status than initially hoped. It is a successful test branch, because it exposed the next structure clearly. It is not the sought nondegenerate derivational branch. The immediate implication for the research program is that the next candidate family must now be chosen elsewhere, most naturally in the constraint-assisted \(U\)-sector or another genuinely new completion beyond the present derivative-mixing ansatz.

\end{document}
